% THETA3-Q11 contribution by /root/theta_032, 21 November 2026.
% GPT-6 Astra Pro; inherited effort label not separately exposed.

\subsection{THETA3-Q11: definite periods give an explicit lower bound for the ball norm}\label{Q11}

This bound replaces the unknown upper bound for the ball Gram
matrix in A5 by local holomorphic evaluation constants.  It uses
the actual finite definite restrictions and a fixed smooth
quotient throughout.  No derivative bound for the local
holomorphic coefficient is needed.

\paragraph{Fixed geometric and period normalizations.}
Let $X$ be a compact smooth ball quotient of complex dimension
$n$, or the finite disjoint union of components of a fixed neat
Shimura level.  Let $\mathcal E\subset H^0(X,\Omega_X^n)$ be a
finite-dimensional space in its specified scalar adelic
realization.  Use the geometric Hermitian norm
\[
 \langle\xi,\zeta\rangle_X
       =i^{n^2}\int_X\zeta\wedge\overline\xi,
 \qquad \|\xi\|_X^2=i^{n^2}\int_X\xi\wedge\overline\xi.
 \tag{Q11a}
\]
It is conjugate linear in the first variable.  If an adelic
Petersson measure is used instead, its specified conversion to
(Q11a) must be included; the argument does not identify these
measures without that conversion.

Take a nonempty finite collection of restrictions $R_{a_j}$ of A5, with finite adelic
$a_j$, totally definite $H=\U(n)$, and the fixed canonical
scalar type.  The same argument applies to any such finite
collection; separation is needed only when positive definiteness
on all of $\mathcal E$ is desired.  Give each $[H]$ probability
Haar measure, choose an orthonormal basis $e_{j,\alpha}$ of its
finite restriction module, and write
\[
 B_{j,\alpha}(\xi)
     =\int_{[H]}(R_{a_j}\xi)(h)\,
                              \overline{e_{j,\alpha}(h)}\,dh .
\]
For representatives of
$H(F)\backslash H(\A_{F,f})/K_H(a_j)$, let $r_{j,t}(\xi)$
be the corresponding scalar values.  Their exact masses are
\[
 m_{j,t}=
    \frac{|\Gamma_{j,t}|^{-1}}
         {\displaystyle\sum_s|\Gamma_{j,s}|^{-1}}>0,
 \qquad
 \Gamma_{j,t}=H(F)\cap h_{j,t}K_H(a_j)h_{j,t}^{-1}.
 \tag{Q11b}
\]
Here the finite stabilizers are viewed in the compact real
group.  The formula follows by decomposing the quotient into
$\Gamma_{j,t}\backslash H(F\otimes_\Q\R)$, first giving the
real group volume one and then normalizing total quotient mass.
A stabilizer character can force a coordinate to vanish; such
coordinates may be retained with value zero.  On a neat level
the relevant finite stabilizers are trivial.  In particular the
masses are concrete rational numbers, although the form values
have not thereby been proved algebraic.

Finite-dimensional Parseval gives the exact positive quantity
\[
 \mathcal P(\xi):=\sum_{j,\alpha}|B_{j,\alpha}(\xi)|^2
          =\sum_{j,t}m_{j,t}|r_{j,t}(\xi)|^2.
 \tag{Q11c}
\]
Flatten the pairs $(j,t)$ into a single index $i$.  Each scalar
value is a linear evaluation in the canonical fibre at a point
$x_i\in X$.  Choose an injective holomorphic chart
$\phi_i:B_{R_i}\subset\C^n\longrightarrow X$ with
$\phi_i(0)=x_i$ and $0<r_i<R_i$.  Write
\[
 \phi_i^*\xi=g_i(z)\,dz_1\wedge\cdots\wedge dz_n,
 \qquad r_i^{\rm ev}(\xi)=\gamma_i g_i(0),
 \tag{Q11d}
\]
where $r_i^{\rm ev}$ denotes the scalar evaluation formerly
indexed by $(j,t)$, so it is distinct from the radius $r_i$.
The nonzero constant $\gamma_i$ is fixed by the actual
coefficient frame: explicitly it is the value of that
evaluation functional on the coordinate canonical generator
at $x_i$.  Thus it includes every prescribed frame or
rigidification factor.  Neither $\gamma_i=1$ nor its
algebraicity is assumed.

Put
\[
 A_i=\frac{2^n\pi^n r_i^{2n}}{n!},\qquad
 c_i=\frac{m_i|\gamma_i|^2}{A_i},\qquad
 C_{\rm eval}=\sum_i c_i.
 \tag{Q11e}
\]
The radii, frame constants and masses are all data for the
chosen quotient; changing its level requires recording the
corresponding changed data.

\paragraph{The explicit bound.}
For every $\xi\in\mathcal E$ one has
\[
 \boxed{\quad
       \mathcal P(\xi)\leq C_{\rm eval}\|\xi\|_X^2,
       \qquad
       \|\xi\|_X^2\geq
                \frac{\mathcal P(\xi)}{C_{\rm eval}}.
       \quad}
 \tag{Q11f}
\]
In particular any explicit lower certificate for the sum of
definite period squares gives a lower certificate for the
geometric ball norm.  This does not use an upper bound for
$G_{\rm ball}$.

\textbf{Proof, including the local constant.}
For a holomorphic function $g$ on $B_R$ and $0<r<R$,
\[
       \int_{B_r}|g(z)|^2\,dV(z)
          \geq \frac{\pi^n r^{2n}}{n!}|g(0)|^2.
 \tag{Q11g}
\]
To see the mean identity directly, fix $w\in B_r$.  The
one-variable function $u\mapsto g(uw)$ is holomorphic on a
disc containing the closed unit disc, so its circle average is
$g(0)$.  Integrate this identity in $w$.  Rotation invariance
of $B_r$ and Fubini show that the average of $g$ on $B_r$ is
$g(0)$.  Cauchy--Schwarz then gives (Q11g), using
$\operatorname{vol}(B_r)=\pi^nr^{2n}/n!$.

In coordinates $z_j=x_j+iy_j$, the order of the exterior
factors gives exactly
\[
 i^{n^2}\,dz_1\wedge\cdots\wedge dz_n\wedge
        d\bar z_1\wedge\cdots\wedge d\bar z_n
        =2^n\,dx_1\,dy_1\cdots dx_n\,dy_n.
\]
Applying (Q11g) to $g_i$ therefore proves
\[
 \int_{\phi_i(B_{r_i})}i^{n^2}\xi\wedge\bar\xi
          \geq A_i |g_i(0)|^2.
 \tag{Q11h}
\]
The chart is injective, so this is an integral over a genuine
subset of $X$, with no covering multiplicity.  Multiplication
by $m_i|\gamma_i|^2/A_i$ gives
\[
 m_i|r_i^{\rm ev}(\xi)|^2
       \leq c_i\int_{\phi_i(B_{r_i})}i^{n^2}\xi\wedge\bar\xi
       \leq c_i\|\xi\|_X^2.
\]
Summing and using (Q11c) proves (Q11f).
\hfill$\square$

The same proof retains overlap information.  If
$U_i=\phi_i(B_{r_i})$, one may replace $C_{\rm eval}$ by
\[
 C_{\rm ov}=\sup_{x\in X}\sum_{i:x\in U_i}c_i
         \leq C_{\rm eval}.
 \tag{Q11i}
\]
Indeed the sum of the local integrals is the integral against
the displayed nonnegative weight function.  In particular,
if no point lies in more than $\kappa$ chosen charts, then
$C_{\rm ov}\leq\kappa\max_i c_i$.
The single-point consequence is also sharper than a derivative
envelope: if $|g_i(0)|\geq b$, then
$\|\xi\|_X^2\geq A_i b^2$ for every recorded injective radius.
It removes the Lipschitz hypothesis from the pointwise estimate
in Q10 and its factor-four loss.

\paragraph{A mixed-period and component-extraction certificate.}
Suppose $\xi,\xi'\in H^0(X,\Omega_X^n)$ are expressed on the
same quotient and in the same canonical line bundle, and
\[
       \|\xi'-\xi\|_X\leq\eta\|\xi\|_X,
       \qquad 0\leq\eta<1.
 \tag{Q11j}
\]
They may have been constructed from different theta types;
that origin does not replace the asserted norm comparison.
Cauchy--Schwarz and (Q11f) give
\[
 \begin{split}
 \operatorname{Re}\langle\xi,\xi'\rangle_X
      &\geq(1-\eta)\|\xi\|_X^2,\\
 |\langle\xi,\xi'\rangle_X|
      &\geq\frac{1-\eta}{C_{\rm eval}}
                    \sum_{j,\alpha}|B_{j,\alpha}(\xi)|^2.
 \end{split}
 \tag{Q11k}
\]
For completeness, subtract $\|\xi\|_X^2$ from the pairing:
the difference is $\langle\xi,\xi'-\xi\rangle_X$, whose
absolute value is at most $\eta\|\xi\|_X^2$.  This proves
the first inequality and hence the second.

If the constructed comparison form has a finite expansion
$\xi'=\sum_\nu b_\nu\xi'_\nu$, set
$S=\sum_\nu|b_\nu|$.  For $\xi\ne0$, (Q11j) implies $S>0$.
Linearity in the second variable and the triangle inequality
then prove that at least one actual component satisfies
\[
 |\langle\xi,\xi'_\nu\rangle_X|
       \geq\frac{1-\eta}{S C_{\rm eval}}
                    \sum_{j,\alpha}|B_{j,\alpha}(\xi)|^2.
 \tag{Q11l}
\]
If both forms are expanded, the same proof uses the product
of their two coefficient $\ell^1$ sums.  The finite tensor and
character-extraction costs in A3 or Q7 can be inserted into
$S$ once their exact expansions are fixed.  No uniform bound
for those costs is silently supplied here.

\paragraph{A determinant certificate without the ball Gram matrix.}
Choose any basis $\xi_1,\ldots,\xi_D$ of $\mathcal E$ and form
the weighted evaluation matrix
\[
 E_{ia}=\sqrt{m_i}\,r_i^{\rm ev}(\xi_a),\qquad
 H_{\rm fin}=E^*E.
\]
For $\xi=\sum_a u_a\xi_a$ one has
$\mathcal P(\xi)=u^*H_{\rm fin}u$.  Suppose a square
$D$-row minor has determinant of modulus at least $\delta>0$
and $\operatorname{tr}H_{\rm fin}\leq T$.  If $D>1$, then
\[
 \|\xi\|_X^2
    \geq\frac{\delta^2}{C_{\rm eval}}
               \left(\frac{D-1}{T}\right)^{D-1}
                         \sum_a|u_a|^2 .
 \tag{Q11m}
\]
To prove this, Cauchy--Binet gives
$\det H_{\rm fin}\geq\delta^2$.  The product of all but
the smallest eigenvalue is at most
$(T/(D-1))^{D-1}$ by the arithmetic--geometric mean inequality.
Thus its smallest eigenvalue is at least the numerator in
(Q11m) before division by $C_{\rm eval}$.  Apply (Q11f).
For $D=1$ the corresponding bound is
$\delta^2|u_1|^2/C_{\rm eval}$.

This is a finite matrix certificate in prescribed coordinates;
an absolute positive number also requires a lower bound for
$\sum_a|u_a|^2$ for the particular input.  The exact mass and
height/denominator arguments in B5 and B7 may provide the
definite data only under their stated arithmetic normalization
hypotheses.  Rationality of the masses alone proves neither
algebraicity of the entries $r_i^{\rm ev}(\xi_a)$ nor a height
bound for their determinants.

\paragraph{Meaning for the comparison problem.}
The new constant $C_{\rm eval}$ depends on local injective
charts, the exact coefficient frames and finite definite
masses.  These data replace an unknown bound for
$G_{\rm ball}$ in this lower-bound argument.  They still must
be supplied to obtain a numerical certificate, as must any
positive definite-period lower bound and the relative error
in (Q11j).  The argument does not compute those numbers from
signature alone, nor identify a single Bessel period with the
whole ball pairing.  It gives an explicit route from the
finite definite periods to the same geometric norm that
controls the mixed theta period.

\framework{This is a comparison of positive Hermitian forms
and specified fibre evaluations.  It fits the polarization
and correspondence discussion in Grothendieck's
\emph{cote44}, scans 7--8, while retaining the distinction
between fibrewise evaluation and the cohomological operations
discussed in \emph{cote16}, scans 43--46.  The analytic estimate
does not turn the evaluation matrix or the frame correction
of A5 into a rational algebraic correspondence.}
